\documentclass[10pt]{article}

\usepackage[T1]{fontenc}
\usepackage{lmodern}
\usepackage[a4paper,margin=19mm]{geometry}
\usepackage{booktabs}
\usepackage{tabularx}
\usepackage{xcolor}
\usepackage{hyperref}

\hypersetup{colorlinks=true,urlcolor=blue,linkcolor=blue,
  pdftitle={Hedging model inventory and AI-agent feedback}}
\setlength{\parindent}{0pt}
\setlength{\parskip}{0.45em}
\renewcommand{\arraystretch}{1.14}

\begin{document}

\begin{center}
  {\LARGE Hedging model inventory and implementation status}\par
  {\large AI-agent feedback, 24 October 2026}
\end{center}

\noindent\fcolorbox{black}{gray!10}{\parbox{0.96\linewidth}{
\textbf{Provenance and status.} This is provisional feedback prepared by AI
coding agents and checked by an independent QuantAgent. It records implementation
observations, not a peer-reviewed result or a claim that unfinished models hedge
well. Measured outcomes, small execution smokes, and agent assessments are kept
separate. The older \texttt{hedging\_benchmark\_summary.tex} concerns a different
611-origin/20-target cohort and must not be merged with the current April--May
short-end experiment.}}

\section{Dataset and experiment}

The only training and evaluation input for this benchmark is the normalized
public-candidate dataset:
\begin{center}
\texttt{spx-2022-2025-hedging-public-candidate-v1}
\end{center}
It has 891 retained dates from 3 January 2022 to 29 August 2025 and
2,014,740 recommended hedge-screen quote-side coordinates. Its manifest
SHA-256 is:
\begin{center}\small
\texttt{2add65d6a1dda4146214479788cdd8e5df652562812510f607a5e8331b313541}
\end{center}
It stores reconstructable SANOS and DLV states/geometries and separate fitted
model price and Black-IV arrays, including ATM volatility by expiry. It also
stores normalized half-spreads, forward/spot ratios, discount factors, and
forward-adjusted spot returns. Fitted SANOS prices are the common reference
marks; they are \emph{not} unavailable market midpoints. The artifact excludes
original WRDS bids, asks, prices, identifiers, volume, open interest, vendor
Greeks, and vendor IVs. Its \texttt{public\_candidate} label is a technical
classification, \emph{not} redistribution permission.

The current short-end generator grid has business-day expiries
$1,2,3,4,5,7,10$ and 19 fixed normalized-moneyness (NMS) nodes from $-3.5$
to $3.5$. Actual forward-relative strikes float by date and expiry using that
expiry's SANOS ATM volatility. The eligible hedge book uses spot and OTM
options inside $(-2,2)$ NMS at expiries through five business days, or the
first later expiry if needed. April and May 2024 each contain 22 completed
daily origins and ten one-day digital targets in the current comparison.
Historical rolling estimates use information available through the decision
date, without future returns; model initialization uses 2022. A newly stated
rolling requirement is that every neural model's date-$t$ weights start from
the immediately previous fitted date, with authenticated checkpoint lineage.
The 2022 initializer is used only at the first rolling origin. This requirement
must be verified before accepting new rolling results.

\section{Models and what they do}

\begin{tabularx}{\linewidth}{@{}p{28mm}X p{29mm}@{}}
\toprule
Method & Construction and hedge & April--May evidence \\
\midrule
Unhedged; BS delta+vega & No hedge; or spot delta plus the two vanilla options
bracketing the digital strike, weighted for vega and re-delta-hedged. & Baselines complete \\
SANOS Greek & DLV local-vol repricing gives spot delta and eight rolling DLV
factor vegas; a minimum-spread-cost factor hedge uses eligible options, then
spot removes residual delta. For a settling 1D digital, future-DLV vega is
zero, so this hedge trades spot only. & Complete \\
SANOS parametric & Uses the same Greeks and a 252-transition spot/factor
covariance to minimize one-day P\&L variance plus executable cost. & 44-date
run complete \\
SANOS statistical & Replays historical SANOS-based hedge and target returns;
optimizes the empirical scenario objective with spread costs and no Gaussian
return assumption. & Corrected 44-date run complete \\
SANOS normal-statistical & Normal approximation using covariance of
SANOS returns and the same cost-aware objective. Legacy DLV covariance runs
are \emph{not} this method. & 44-date run complete \\
GROMM Greek & Proposed Q expectation and spot plus eight rolling DLV-factor
Greeks from the GROMM stochastic-DLV simulator. & No result yet \\
GROMM parametric & Proposed covariance-weighted GROMM spot/factor Greek hedge.
& No result yet \\
GROMM statistical & Proposed optimization on GROMM simulated P or Q
instrument/target scenarios, with the declared measure identified. & No result yet \\
GROMM normal-statistical & Proposed Gaussian approximation to those GROMM
scenario returns. & No result yet \\
VolGAN & Conditional GAN jointly generates next-day normalized spot returns
and implied-volatility changes under historical P. The original full surface
and a separately labelled eight-PC adaptation are retained. & Both 2022
initializers trained; PCA-8 rolling cell complete \\
AD--Seq--Vol; FT & Conditional sequential diffusion generates spot/IV paths
under P. FT adds post-training static-arbitrage penalties to the frozen base.
& Base trained; FT running \\
FuNVol & Functional-basis/PCA reduction of IV surfaces with a joint spot/factor
neural SDE under P. The present implementation is a single-asset research
adaptation of the published multi-asset model. & Rolling gate under review \\
\bottomrule
\end{tabularx}

\section{Measured April--May evidence}

The table reports equal-weight mean across ten one-day digital targets of
each target's mean absolute \emph{pure-hedge} error, in inception-spot units.
These are 22-date monthly results with the common SANOS reference mark and
realized accounting. The legacy statistical and DLV-replay covariance cells
shown are mean--variance with risk parameter 10 and sample demeaning; this is
not the full objective/risk tensor. The DLV-replay rows are completed auxiliary methods,
not the proposed SANOS normal-statistical method. VolGAN uses the same
mean--variance risk parameter and demeaning, but with 10,000 generated paths
per origin and causally linked daily checkpoints. It is the exploratory
eight-PC adaptation, not the original full-surface implementation.

\begin{center}
\begin{tabular}{@{}llrr@{}}
\toprule
Model & Hedge method & April 2024 & May 2024 \\
\midrule
Baseline & Unhedged & 0.011128 & 0.010825 \\
Baseline & Black delta+vega & 0.006416 & 0.005733 \\
SANOS & Greek & 0.007284 & 0.007479 \\
SANOS & Parametric & 0.010585 & 0.008980 \\
SANOS & Statistical & 0.003707 & 0.005035 \\
SANOS & Normal-statistical & 0.007735 & 0.005843 \\
Legacy DLV & Statistical replay & 0.003706 & 0.005034 \\
Legacy DLV & Covariance & 0.007889 & 0.005903 \\
Legacy DLV & Covariance, PCA-8 & 0.007892 & 0.005940 \\
VolGAN (PCA-8) & Scenario hedge & 0.011301 & 0.010843 \\
\bottomrule
\end{tabular}
\end{center}

The earlier daily ``Historical-SANOS'' implementation actually replayed DLV,
not SANOS, model states. Its measured result above is retained only as a
clearly labelled legacy diagnostic; the corrected SANOS-state replay completed
all 440 target-date cells in a separate, authenticated cohort, with no solver
nonconvergence or cap hit. The two replays produce similar, but not identical,
results. The SANOS normal-statistical fit also completed all 440 cells with no
solver nonconvergence or cap hit. The SANOS parametric fit likewise completed
440 cells with no solver nonconvergence or cap hit, but its hedge error exceeds
the Greek hedge in both months. SANOS Greek is the method
previously named SANOS--LV, not a different model.
For a settling 1D digital, both the earlier and current Greek runs have zero
future-DLV factor vega and trade spot only; statistical replay may also trade
options. This difference in hedge instruments does \emph{not} explain a change
in ranking between the two runs. The earlier 283-date comparison used a
different period, an entropy-$\lambda=1$ historical policy rather than the
present mean--variance risk-10 policy, and a $0.01\%$ rather than $1\%$
relative-spot finite-difference bump for the Greek delta. On the matched
April dates, the earlier Greek and statistical mean absolute errors were
$0.007738$ and $0.007388$; the current errors are $0.007284$ and $0.003706$.
On the matched May dates, the corresponding pairs were $0.006855$, $0.007571$
and $0.007479$, $0.005034$. The ranking reversal is therefore not a model
rename or evidence of a new Greek hedge. VolGAN has 440 completed target-date results, zero
numerical nonconvergences and zero cap hits, but its hedge error is marginally
above unhedged and its conditional shocks are under-dispersed. AD--Seq--Vol,
its FT variant, FuNVol, and GROMM variants have no validated
April--May hedge result yet. An absent result is not a zero or an omitted
favourable measurement.

\section{Implementation notes and difficulties}

\textbf{VolGAN.} The published-length full-surface adaptation trained for
10,000 CPU epochs on 213 one-session examples. It generated too little
spot-return variation when the condition was held fixed. Eight PCs explain
99.21\% of training and 97.35\% of validation log-IV \emph{levels}, and 96.65\%
and 91.54\% of log-IV \emph{increments}. The separately labelled eight-PC
variant also completed 10,000 epochs. Across 20 validation states it generated
0.515\% within-state daily spot-return standard deviation, against 1.435\%
out-of-sample residual standard deviation from a predictor fitted only on the
training block. Generated IV-change RMS was 0.0411 versus 0.1228 for those
residuals in log-IV units. Its final validation loss was 1.047 versus best
observed 0.896; the initializer checkpoint stores final rather than best
weights. The 44-origin April--May eight-PC rolling hedge has completed with
10,000 paths per origin, but remains \emph{exploratory}, not evidence that
VolGAN is conditionally calibrated. Compressing the IV output improved but did
not cure conditional under-dispersion.

\textbf{AD--Seq--Vol and FT.} The base completed 3,000 CPU epochs, with best
validation loss 0.047693. FT is training on CPU from the SHA-checked base checkpoint in a
separate directory; no FT or rolling hedge outcome is yet reported here.

\textbf{FuNVol.} The CPU implementation, fixed-contract repricer, 2022
initializer, and causal 2023--March rolling predecessor chain are in place.
The 1 April 2024 generated one-day spot standard deviation was $0.3327\%$,
versus $0.7777\%$ in the trailing historical comparison, below the
precommitted minimum ratio of $0.5$. No April--May hedge record was published;
the conditional-state versus model-underdispersion cause is under review.
Positive IV alone is not a static- or dynamic-arbitrage guarantee.

\textbf{Shared limitations.} Generated spot/IV models currently hold
forward/spot ratios and discount factors fixed within simulated transitions.
On a 40-date diagnostic, the median same-strike fixed-grid IV approximation
error in digital marks was 0.000711, versus 0.006461 from using a frozen
rather than realized next-day forward. This is a model-law issue, not merely
grid resolution. P-scenario generators are not automatically Q pricing models;
measure claims and drift removal must remain explicit. Rolling optimization
also needs true previous-date checkpoint inheritance, solver/cap diagnostics,
and identical executable spread costs before comparison.

\section{AI-agent assessment and next checks}

\textbf{AI-agent assessment, not a measured result:} the tiny 2022 sample
relative to a high-dimensional spot/IV output is a plausible contributor to
VolGAN's conditional under-dispersion, but PCA coverage alone cannot establish
the cause. The next diagnostic should compare conditional spot and IV residual
scales, cross-dependence, static-arbitrage rates, and held-out digital marks on
each dated checkpoint before interpreting hedge errors.

\textbf{QuantAgent review:} the older generated summary concerns a distinct
full-period experiment. Keep its table and the current April--May short-end
table separate; never merge counts or rankings. The April--May rolling jobs
must also prove previous-date weight lineage rather than repeatedly loading
the 2022 initializer. These are verification requirements, not completed
performance findings.

\vspace{0.5em}
\textit{Primary references:} Vuleti\'c--Cont, \href{https://doi.org/10.1080/1350486X.2025.2471317}{VolGAN};
Han et al., \href{https://arxiv.org/abs/2609.13402}{AD--Seq--Vol};
Choudhary--Jaimungal--Bergeron, \href{https://doi.org/10.1080/14697688.2024.2396977}{FuNVol}.

\end{document}
